% LaTeX for Distribution Shift Analysis
% Experiment 02: Feature Distribution Shift (Statistical Analysis)

\subsection{Distribution Shift Analysis}
\label{sec:distribution_shift}

BanditGPT ships with warm-start priors learned from 80K RouteLLM battles that provide immediate routing intelligence on Day~1 (Table~\ref{tab:performance_gap}). But a natural question arises: \emph{will these priors still be useful when the library encounters real user traffic?} If the deployment distribution differs from the training data---a phenomenon known as \emph{covariate shift}~\cite{shimodaira2000improving}---the shipped priors may degrade the very Day~1 advantage they were designed to provide.

This section quantifies the shift between RouteLLM warmup data and LMSYS evaluation prompts, establishing the practical motivation for two library capabilities: (1)~online adaptation via LinUCB, and (2)~Corralling as bounded-cost insurance against more severe mismatch (evaluated in Table~\ref{tab:performance_gap}).

\subsubsection{Methodology}

We employ the Population Stability Index (PSI)~\cite{yurdakul2018statistical} to quantify distribution mismatch. Critically, we extract features using the \emph{production} \texttt{FeatureService.extract\_features()} code path---the same pipeline that processes every prompt in deployment. This ensures the measured shift is exactly what the router's LinUCB agent observes, not a comparison in an abstract feature space. PSI is computed along PC1 using quantile-based bins ($B=20$) from the 80K RouteLLM reference distribution:

\begin{equation}
\text{PSI} = \sum_{i=1}^{B} \left( (\%P_i - \%Q_i) \cdot \ln\left(\frac{\%P_i}{\%Q_i}\right) \right)
\label{eq:psi}
\end{equation}

\noindent where $\%P_i$ and $\%Q_i$ are the proportions in bin $i$ for the reference ($P$, warmup) and deployment ($Q$, LMSYS evaluation) distributions, respectively. Standard thresholds: PSI $< 0.1$ (no significant shift), $0.1$--$0.2$ (moderate), $0.2$--$0.25$ (significant), $\geq 0.25$ (substantial).

\subsubsection{Results}

Our analysis yields \textbf{PSI = 0.225 (95\% CI: [0.194, 0.285])}, placing the shift in the ``significant'' range. In practical terms, this is a \textbf{small but reliable effect}: Cohen's d $= 0.33$ (small), KS test D $= 0.119$ ($p < 10^{-22}$), mean PC1 shift $= +0.060$. We emphasize effect sizes over p-values: with $N > 10{,}000$ samples, even trivial differences reach significance. The practically relevant finding is that deployment prompts occupy a measurably different region of feature space than the warmup prior.

The shift is concentrated on the leading PCs (where between-domain variation lives) and attenuates at higher PCs: PC1 PSI $= 0.225$, PC1--5 avg $= 0.197$, all 33 features avg $= 0.089$. See Appendix~\ref{app:distribution_shift} for the full robustness analysis.

\begin{figure}[t]
\centering
\includegraphics[width=0.95\textwidth]{experiments_v1/02_figure/results/figure2_distribution_shift.png}
\caption{\textbf{Distribution shift between shipped priors and deployment data.}
All features extracted via the production \texttt{FeatureService} pipeline.
\textbf{Top:} Warmup priors (Blue, 80K RouteLLM battles) and deployment prompts (Red, LMSYS evaluation) occupy measurably different regions of feature space: PSI $= 0.225$ (95\% CI: $[0.194, 0.285]$), Cohen's d $= 0.33$. The priors provide a strong but imperfect initialization.
\textbf{Bottom:} Warmup data decomposed by ground truth reward gaps. Mixtral-Sufficient tasks (Gap $\leq 0$; 32\%) and GPT-4-Turbo-Required tasks (Gap $= 1$; 68\%) are separated along PC1 (Cohen's d $= 0.23$, $>$99\% overlap). This weak per-prompt signal is exactly the type contextual bandits exploit by accumulating evidence over many rounds.}
\label{fig:distribution_shift}
\end{figure}

\subsubsection{Task Difficulty Structure}

Figure~\ref{fig:distribution_shift} (bottom) decomposes the warmup prior by ground truth model performance. We partition by reward gap ($R_{\text{GPT-4-Turbo}} - R_{\text{Mixtral}}$):
\begin{itemize}[leftmargin=*,noitemsep,topsep=2pt]
    \item \textbf{Mixtral-Sufficient} (Gap $\leq 0$; 32\%): centered at PC1 $= 0.022$
    \item \textbf{GPT-4-Turbo-Required} (Gap $= 1$; 68\%): centered at PC1 $= -0.016$
\end{itemize}

\noindent These categories overlap extensively ($>$99\% range overlap; centroids differ by only 0.038) but are statistically separable (Cohen's d $= 0.23$, Welch's t $= 10.60$, $p < 10^{-25}$). This characterizes the routing problem realistically: features provide a \emph{noisy, weak} signal, not a clean decision boundary. LinUCB does not need clean separation---it needs a \emph{consistent} signal that it accumulates over many rounds into reliable routing decisions.

\subsubsection{Implications}

This analysis answers a question critical for library users: \emph{do shipped priors remain valid on new data?} The answer is nuanced:

\begin{enumerate}[leftmargin=*,noitemsep,topsep=2pt]
    \item \textbf{Priors are a strong but imperfect initialization.} The shift is significant but practically small (d $= 0.33$). Table~\ref{tab:performance_gap} confirms that warmup priors remain beneficial under this shift (28\% lower regret than tabula rasa), but larger shifts in real deployments could degrade prior quality more substantially.
    \item \textbf{Online adaptation provides continuous correction.} LinUCB converts weak per-prompt signals (d $= 0.23$) into effective long-run policies. If shift arises even within the relatively homogeneous LLM evaluation domain, it will be larger where user populations vary more widely~\cite{lu2018learning}.
    \item \textbf{Corralling provides bounded-cost insurance---and the payoff is demonstrated empirically.} Under this moderate shift, Corralling adds only 1--4\% overhead vs tabula rasa (Table~\ref{tab:performance_gap})---a small premium when priors are already well-matched. That premium pays off under more severe conditions tested later in this paper: when a model catastrophically fails, Corralling detects and fails over within 3 steps with 100\% success rate (Figure~\ref{fig:corralling}); when a new model is released, semantic transfer eliminates the cold-start penalty entirely, yielding +39\% improvement over cold start (Figure~\ref{fig:zero-shot-readiness}).
\end{enumerate}

\noindent \textbf{Narrative arc.} Figure~\ref{fig:lmsys_holdout_structure} establishes that PCA features generalize to unseen data. This section shows generalization is \emph{imperfect}. Table~\ref{tab:performance_gap} completes the immediate arc: warmup priors remain beneficial despite the shift, and Corralling's 1--4\% insurance premium is modest. The full return on that investment is demonstrated when conditions deteriorate further: catastrophic model failure triggers 3-step automatic failover (Figure~\ref{fig:corralling}), new model releases are absorbed without cold-start penalty (+39\% vs na\"ive adoption; Figure~\ref{fig:zero-shot-readiness}), and performance is robust across a $20\times$ hyperparameter range (Figure~\ref{fig:sensitivity}). Together: \emph{ship strong priors, adapt online---and when the unexpected happens, Corralling pays for itself}.

Detailed robustness analysis, PCA justification, configuration guidance, and comparison to related work appear in Appendix~\ref{app:distribution_shift}.
